\section{Exact squarefree-support and parity splitting}
\label{sec:parity-split}

Let \(\lambda(n)=(-1)^{\Omega(n)}\) be the Liouville function.  For
every \(n\in\N\),
\begin{equation}
 \boxed{\mu(n)=\mu^2(n)\lambda(n).}
\label{eq:mu-squarefree-lambda}
\end{equation}
Indeed, both sides vanish when \(n\) is not squarefree, while on a
squarefree integer both equal \((-1)^{\omega(n)}\).

\begin{theorem}[Exact two-form split]
\label{thm:two-form-split}
Let \(I\subset\Z\) be a finite interval on which
\(L_1(r),L_2(r)>0\), and let \(W:I\to\C\) be arbitrary.  Then
\begin{equation}
 \boxed{
 \sum_{r\in I}W(r)\mu(L_1(r))\mu(L_2(r))
 =
 \sum_{r\in I}
 W(r)\Qsf(r)\lambda(L_1(r))\lambda(L_2(r)).}
\label{eq:weighted-parity-split}
\end{equation}
The factor \(\Qsf\) is determined entirely by the square-power local
sets \(\Rset_{p,2}\); the remaining factor is a two-point Liouville
sign.
\end{theorem}

\begin{proof}
Apply \eqref{eq:mu-squarefree-lambda} to both forms, multiply, and
sum.
\end{proof}

\begin{remark}[Support is not sign cancellation]
\label{rem:support-not-sign}
The positive density and exact root counts of
\cref{sec:local-root-counts,sec:squarefree-support} concern only
\(\Qsf\).  They do not estimate the signed sum on the right of
\eqref{eq:weighted-parity-split}.  Replacing the Liouville product by
the constant sign \(1\) preserves every local-support statement and
produces no cancellation for nonnegative \(W\).  This is a logical
separation, not a model for the actual Liouville function.
\end{remark}

\subsection{Exact square-divisor expansion}

The identity
\begin{equation}
 \mu^2(n)=\sum_{q^2\mid n}\mu(q)
\label{eq:squarefree-divisor-identity}
\end{equation}
turns the local support into congruence classes.  The determinant
gives a particularly sharp compatibility condition.

\begin{theorem}[Primitive square-divisor compatibility]
\label{thm:square-divisor-compatibility}
Let \(q_1,q_2\) be positive squarefree integers and put
\(g=\gcd(q_1,q_2)\).  For a primitive pair
\((L_1,L_2)\), the system
\begin{equation}
 q_1^2\mid L_1(r),\qquad
 q_2^2\mid L_2(r)
\label{eq:square-divisor-system}
\end{equation}
has an integral solution \(r\) if and only if
\begin{equation}
 \boxed{
 \gcd(q_1,s)=1,\qquad
 \gcd(q_2,a)=1,\qquad
 g^2\mid h_0.}
\label{eq:square-divisor-compatibility}
\end{equation}
When soluble, \eqref{eq:square-divisor-system} defines exactly one
residue class modulo
\begin{equation}
 \lcm(q_1^2,q_2^2)
 =\frac{q_1^2q_2^2}{g^2}.
\label{eq:square-divisor-modulus}
\end{equation}
\end{theorem}

\begin{proof}
The first congruence in \eqref{eq:square-divisor-system} is soluble
if and only if \((q_1,s)=1\).  To see this prime by prime, if
\(p\mid(q_1,s)\), primitivity gives \(p\nmid d\), so there is no root
even modulo \(p\); if \((q_1,s)=1\), there is one root modulo
\(q_1^2\).  The second congruence is analogous.

Under these two coprimality conditions, the individual congruences
specify one class modulo \(q_1^2\) and one class modulo \(q_2^2\).
The generalized Chinese remainder theorem says they are compatible
exactly when their reductions agree modulo
\(\gcd(q_1^2,q_2^2)=g^2\).  Since \(q_1,q_2\) are squarefree, this can
be checked at every prime \(p\mid g\).  Both slopes are then units
modulo \(p\), and \cref{lem:common-root-criterion} with \(k=2\) says
that the roots agree modulo \(p^2\) exactly when \(p^2\mid h_0\).
Thus compatibility is equivalent to \(g^2\mid h_0\).  The unique
class modulo \(\lcm(q_1^2,q_2^2)\) is the standard conclusion of the
generalized Chinese remainder theorem.
\end{proof}

\begin{corollary}[Exact weighted divisor expansion]
\label{cor:exact-weighted-divisor-expansion}
Under the hypotheses of \cref{thm:two-form-split},
\begin{align}
 &\sum_{r\in I}W(r)\mu(L_1(r))\mu(L_2(r))
\notag\\
 &\quad=
 \sum_{\substack{q_1,q_2\ge1\\
                 (q_1,s)=(q_2,a)=1\\
                 \gcd(q_1,q_2)^2\mid h_0}}
 \mu(q_1)\mu(q_2)
 \sum_{\substack{r\in I\\
                 q_1^2\mid L_1(r)\\
                 q_2^2\mid L_2(r)}}
 W(r)\lambda(L_1(r))\lambda(L_2(r)).
\label{eq:exact-weighted-divisor-expansion}
\end{align}
The outer sum is finite: only \(q_i^2\le
\max_{r\in I}L_i(r)\) can contribute.  Each nonempty inner sum is
over one residue class modulo
\(\lcm(q_1^2,q_2^2)\).
\end{corollary}

\begin{proof}
Insert \eqref{eq:squarefree-divisor-identity} twice into the
right-hand side of \eqref{eq:weighted-parity-split}, interchange the
finite sums, and apply
\cref{thm:square-divisor-compatibility}.  Terms with nonsquarefree
\(q_i\) vanish because \(\mu(q_i)=0\).
\end{proof}

\begin{remark}[What the divisor expansion exposes]
\label{rem:divisor-expansion-exposes}
The overlap \(\gcd(q_1,q_2)\) is supported only on prime squares
dividing \(h_0\), a fixed finite resource.  However,
\eqref{eq:exact-weighted-divisor-expansion} still contains Liouville
correlations on progressions whose moduli, endpoints, and
coefficients can grow.  The exact local resolution therefore
organizes the parity problem; it does not eliminate it.
\end{remark}

\subsection{A fixed-form consequence and its boundary}

\begin{corollary}[Unsigned support mass]
\label{cor:unsigned-support-mass}
For a fixed primitive pair,
\[
 \sum_{r\le R}\Qsf(r)
 =
 \Mloc(L_1,L_2)R+o(R).
\]
Consequently,
\[
 \left|
 \sum_{\substack{1\le r\le R\\L_1(r)L_2(r)\ne0}}
 \mu(|L_1(r)|)\mu(|L_2(r)|)
 \right|
 \le
 \Mloc(L_1,L_2)R+o(R).
\]
\end{corollary}

\begin{proof}
The first statement is
\cref{thm:fixed-pair-squarefree-density}.  For every nonzero integer
\(n\), the identity
\(\mu(|n|)=\mu^2(|n|)\lambda(|n|)\) holds.  Apply it to both forms
and use \(|\lambda|=1\); the at most two integral zeros have already
been excluded.
\end{proof}

\begin{remark}[No nontrivial signed saving]
The last upper bound is the support-mass bound and is generally of
order \(R\).  It is not a Chowla estimate and supplies no
\(o(R)\) cancellation.
\end{remark}
